\documentclass[12pt]{article}

% DEFAULT PACKAGE SETUP
\usepackage{setspace,graphicx,epstopdf,amsmath,amsfonts,amssymb,amsthm,versionPO,threeparttable}
\usepackage{marginnote,datetime,enumitem,subfig,rotating,fancyvrb,tabularx,booktabs}
\usepackage{float}
% \usepackage[longnamesfirst]{natbib}
\usepackage[]{natbib}
\newcolumntype{Y}{>{\raggedleft\arraybackslash}X}% raggedleft column X
\interfootnotelinepenalty=10000 %% Completely prevent breaking of footnotes

% Define \numberthis
\newcommand{\numberthis}{\addtocounter{equation}{1}\tag{\theequation}}

\usdate
\usepackage{lscape}
% \usepackage{endfloat}
% \usepackage[nolists]{endfloat}
\usepackage{rotating} 

\usepackage[justification=justified]{caption}
% \usepackage{subcaption}
% These next lines allow including or excluding different versions of text
% using versionPO.sty

\excludeversion{notes}  % Include notes?
\includeversion{links}  % Turn hyperlinks on?

\usepackage{eurosym}
% Turn off hyperlinking if links is excluded
\iflinks{}{\hypersetup{draft=true}}

% Notes options
\ifnotes{%
 \usepackage[margin=1in,paperwidth=10in,right=2.5in]{geometry}%
 \usepackage[textwidth=1.4in,shadow,colorinlistoftodos]{todonotes}%
}{%
 \usepackage[margin=1in]{geometry}%
 \usepackage[disable]{todonotes}%
}

\DeclareMathOperator*{\argmin}{arg\,min}


% Allow todonotes inside footnotes without blowing up LaTeX
% Next command works but now notes can overlap. Instead, we'll define
% a special footnote note command that performs this redefinition.
% \renewcommand{\marginpar}{\marginnote}%

% Save original definition of \marginpar
\let\oldmarginpar\marginpar

% Workaround for todonotes problem with natbib (To Do list title comes out wrong)
\makeatletter\let\chapter\@undefined\makeatother % Undefine \chapter for todonotes

% Define note commands
\newcommand{\smalltodo}[2][] {\todo[caption={#2}, size=\scriptsize, fancyline, #1] {\begin{spacing}{.5}#2\end{spacing}}}
\newcommand{\rhs}[2][]{\smalltodo[color=green!30,#1]{{\bf RS:} #2}}
\newcommand{\rhsnolist}[2][]{\smalltodo[nolist,color=green!30,#1]{{\bf RS:} #2}}
\newcommand{\rhsfn}[2][]{% To be used in footnotes (and in floats)
 \renewcommand{\marginpar}{\marginnote}%
 \smalltodo[color=green!30,#1]{{\bf RS:} #2}%
 \renewcommand{\marginpar}{\oldmarginpar}}
% \newcommand{\textnote}[1]{\ifnotes{{\noindent\color{red}#1}}{}}
\newcommand{\textnote}[1]{\ifnotes{{\colorbox{yellow}{{\color{red}#1}}}}{}}

% Command to start a new page, starting on odd-numbered page if twoside option
% is selected above
\newcommand{\clearRHS}{\clearpage\thispagestyle{empty}\cleardoublepage\thispagestyle{plain}}

% Number paragraphs and subparagraphs and include them in TOC
\setcounter{tocdepth}{2}

% JF-specific includes:

\usepackage{indentfirst} % Indent first sentence of a new section.
\usepackage{endnotes} % Use endnotes instead of footnotes
\usepackage{rfs}  % JF-specific formatting of sections, etc.
% \usepackage[labelfont=bf,labelsep=period]{caption} % Format figure captions
\captionsetup[table]{labelsep=none}

\usepackage{color}%Color
% Define theorem-like commands and a few random function names.
\newtheorem{condition}{CONDITION}
\newtheorem{corollary}{COROLLARY}
\newtheorem{proposition}{PROPOSITION}
\newtheorem{obs}{OBSERVATION}
\newtheorem{lem}{LEMMA}
\newtheorem{theorem}{THEOREM}
\newtheorem{assumption}{ASSUMPTION}
\newcommand{\argmax}{\mathop{\rm arg\,max}}
\newcommand{\sign}{\mathop{\rm sign}}
\newcommand{\defeq}{\stackrel{\rm def}{=}}

\usepackage{multirow}
\usepackage{booktabs}


\usepackage{amsmath}
\usepackage{amsthm}
\usepackage{amsfonts}
\usepackage{relsize}
\usepackage{tikz}
\usepackage{pgfplotstable} 
\usepgfplotslibrary{dateplot}
\usetikzlibrary{fit}
\usepackage{graphicx}
\usepackage{eurosym}
\usepackage{standalone}
% \usepackage{filecontents}

\usepackage[section]{placeins} % keep figures in the same section


\usepackage{appendix}
\appendixtitleon
\newcommand\independent{\protect\mathpalette{\protect\independenT}{\perp}}
\def\independenT#1#2{\mathrel{\rlap{$#1#2$}\mkern2mu{#1#2}}}

%\onehalfspacing
%\linespread{1.25}

\setstretch{1.5}
%\doublespacing

\definecolor{zahi}{RGB}{255,0,0}
\newcommand{\zahi}[1]{\textcolor{zahi}{[zahi: #1]}}



%%%%% Pgfplots %%%%%
\usepackage{pgfplots}
\usepgfplotslibrary{groupplots}

\pgfplotsset{compat=newest}
\usetikzlibrary{patterns}

%%%% Table caption package %%%%%
\usepackage{ragged2e}
\usepackage{longtable}

%%%% Figure imports %%%%
\usepackage{import}

%%%% Color cells in tables %%%%
\usepackage{tabulary,colortbl}
\newcolumntype{C}[1]{>{\centering\let\newline\\\arraybackslash\hspace{0pt}}m{#1}}


\definecolor{andrea}{RGB}{255,0,0}
\newcommand{\andrea}[1]{\textcolor{andrea}{[andrea: #1]}}

\definecolor{yang}{RGB}{255,0,0}
\newcommand{\yang}[1]{\textcolor{yang}{[yang: #1]}}
\newcommand{\jay}[1]{\textcolor{yang}{[j: #1]}}


\usepackage[colorlinks=true,linkcolor=red,anchorcolor=black,citecolor=black, filecolor=black,menucolor=black,runcolor=black,urlcolor=black]{hyperref}
\usepackage{parskip}



%%%%%%%%%%%%%%%%% For plots: get min, max etc


\newcommand*{\GetElement}[4]{%
    \pgfplotstablegetelem{#2}{#3}\of{#1}%
    \let#4\pgfplotsretval
}

% get the first element of a column and store it in a macro
% #1: table
% #2: column (name or [index]<index>)
% #3: macro for value
\newcommand*{\GetFirstElement}[3]{%
    \GetElement{#1}{0}{#2}{#3}%
}

% get the last element of a column and store it in a macro
% #1: table
% #2: column (name or [index]<index>)
% #3: macro for value
\newcommand*{\GetLastElement}[3]{%
    \pgfplotstablegetrowsof{#1}%
    \pgfmathsetmacro\lastrowidx{int(\pgfplotsretval - 1)}%
    \GetElement{#1}{\lastrowidx}{#2}{#3}%
}

% get maximum and minimum value of a column and store them in macros
% #1: table
% #2: column
% #3: macro for max
% #4: macro for min
\newcommand*{\GetMinMax}[4]{%
    \pgfmathsetmacro#3{-16383}%
    \pgfmathsetmacro#4{16383}%
    \pgfplotstableforeachcolumnelement{#2}\of{#1}\as\cellValue{%
        \ifx\cellValue\@empty\else
            \pgfmathsetmacro{#3}{max(#3,\cellValue)}%
            \pgfmathsetmacro{#4}{min(#4,\cellValue)}%
        \fi
    }
}

% get average of column and store it in a macro
% #1: table
% #2: column
% #3: macro for average
\newcommand*{\GetAverage}[3]{%
    \pgfplotstablegetrowsof{#1}%
    \let\rownumber\pgfplotsretval
    \def#3{0}%
    \pgfplotstableforeachcolumnelement{#2}\of{#1}\as\cellValue{%
        \ifx\cellValue\@empty
            % don't count emplty cells
            \pgfmathsetmacro{\rownumber}{\rownumber - 1}%
        \else
            \pgfmathsetmacro{#3}{#3 + \cellValue}%
        \fi
    }
    \pgfmathsetmacro{#3}{#3 / \rownumber}%
}

%%%%%%%%%%%%%%%%%%%%%%%%%%%%%%%%%%%%%%%%%%%%%%%%%%%%%%

%%%% Quintile colors (Green to Red) %%%% 
\definecolor{q1}{rgb}{ .973,  .412,  .42}
\definecolor{q2}{rgb}{ .992,  .749,  .486}
\definecolor{q3}{rgb}{ 1,  .922,  .518}
\definecolor{q4}{rgb}{ .678,  .827,  .498}
\definecolor{q5}{rgb}{ .388,  .745,  .482}
%%%%%%%%%%%%%%%%%%%%%%%%%%%%%%%%%%%%%%%%

% optional: end floats
\usepackage[nolists]{endfloat}

\begin{document}
\setlist{noitemsep}  % Reduce space between list items (itemize, enumerate, etc.)
\onehalfspacing      % Use 1.5 spacing
\renewcommand{\thempfootnote}{\fnsymbol{mpfootnote}}


\vspace{1.5cm}
 \author{
% {\large Jiacui Li\hspace{0.00in}}
 }

\title{\vspace*{0.0in}\Large \bf 
Government Intervention in Stock Markets: the Role of Expectations
% \thanks{\rm  Li (\url{jiacui.li@eccles.utah.edu}) is at David Eccles School of Business, University of Utah. }
}

% \vspace{0.2cm}

\date{\today}              % No date for final submission

% Create title page with no page number

\maketitle


\thispagestyle{empty}
\onehalfspacing
\centerline{\bf ABSTRACT}

\noindent

\medskip

\vspace{0.1cm}

% \noindent \textit{JEL classification:} G11, G12, G14\\
% \noindent \textit{Keywords:} Lead-Lag, Cross-Stock Momentum, Factor Momentum, Reversal, Underreaction

\thispagestyle{empty}
\setcounter{page}{1}


%%%%%%%%%%%%%%%%%%%%%
\newpage
\section{Introduction}

% Motivation
Stock market is often seen as a place for ``price discovery,'' and government intervention to stabilize markets is relatively rare. Such interventions have mostly occurred in crisis periods and through relatively indirect tools (such as short-selling bans, central bank lending facility operations, etc.), and they are typically reserved for extreme circumstances. People are often skeptical about both the desirability and feasibility of such interventions. At a minimum, they are uncommon.

% Motivation, continued
However, things are changing. In the world’s second-largest economy, China, the national team (henceforth NT) intervened in ``crisis mode'' in 2015, but has since begun to regularize/institutionalize such interventions. As far as we can tell, this approach is here to stay. Authorities have recently called for a regular stock market stabilization scheme and asked the NT to be ``more regular'' in what it does. As such, a natural question is: can such interventions be effective? Or, what is required for them to be effective?

% The question
We start from the observation that successful intervention is not easy. Stock markets are often huge and far exceed the financial capacity of the NT as a participant. In China, for example, the stock market’s capitalization is enormous. Even at the peak of intervention, NT holdings were never more than a few percent of the market, so it cannot single-handedly ``pin down'' prices. Second, the initial 2015 intervention did not immediately succeed, as we document later---some of it backfired and sowed confusion. Therefore, whether and how such intervention can be effective is a worthwhile research question in itself.

% Model
To guide thinking about this question, we start with a simple two-period model. NT believes that investors are overly pessimistic about the stock market relative to true fundamentals, and the stock market is in free fall. (The opposite case of cooling an overheating market is symmetric.) NT is modeled as a risk-averse agent (and thus does not have infinite capacity) that wants to stabilize markets. It can decide how much to buy in both rounds. It can also decide whether to 1) release information about its purchases, and 2) pre-commit to further conditional buying. Investors are rational, except for their overly pessimistic belief about fundamentals. Further, NT has a type $\theta \in \{0,1\}$: $\theta = 0$ means NT will not intervene, which can arise from NT not caring enough about the stock market or not having the tools to intervene; $\theta = 1$ means NT can intervene. Investors do not know the type and have a prior with $P(\theta) < 1$.\footnote{This way, the very fact that NT announced that ``oh I have intervened'' is big news in itself, as it moves the belief to $P(\theta) = 1$.} We now discuss the predictions one by one and connect them to the empirics.

% Predictions
The core insight from the model is that NT's behavior has both 1) direct price effects via buying and 2) indirect effects via changing expectations. Changes in expectations arise from two channels. First, knowing that NT has intervened updates beliefs about $\theta$. If investors know that NT cares and can intervene going forward, they revise upward their belief about the second-period stock price, and thus revise their demand upward. Second, if NT further pre-commits to its behavior in period 2, investors become even more optimistic today. Both channels make NT's intervention more effective.

We find several layers of empirical evidence for this view. First, we document that NT intervention has clear ``announcement effects,'' not just ``direct effects.'' We show this in both market-level event studies and better-identified firm-level announcements. At the market level, there were \textcolor{blue}{a couple of episodes} where announcements affected price levels. To get better identification, we further examine stock-level disclosures of NT behavior. There are periodic announcements of NT stock-level positions via end-of-quarter firm disclosures, ``crossing the 5\% threshold and thus having to announce,'' etc., and we conduct event studies around them. These event studies show clear price effects that are distinct from direct price-pressure effects (because the transactions have already been completed). These findings show that indirect effects are real.\footnote{J to all: the benefit of these stock-level studies is that they are well-identified: people will believe that disclosure has a causal effect. The drawbacks are: 1) these are stock-level, rather than market-level, interventions; 2) people may question whether NT have a CHOICE of whether to disclose in each case and thus this is not an ``exogenous shock'' (randomly choosing whether to disclose or not), but none of those can defeat that we prove ``NT disclosure has price effect.''}

We then go one step deeper into the mechanism. In the model, disclosure of NT intervention changes investor \textit{demand}. Our model has a representative agent and thus no trading to speak of. In a model with heterogeneous agents, we would observe trading flows, and we indeed find that investors buy in response to observed (or even inferred) NT trading. First, we document this using TAQ trading flows. (\textcolor{blue}{J to all\footnote{I am still a bit fuzzy on the nature of TAQ order flows in China. For each buyer there is a seller. Thus, is it really possible to say that ``all investors'' adjusted up their demand in response to NT info release?}}) Second, this is not only about realized/announced NT trades; investors also actively try to ``sniff out'' potential NT trades. In 2015, we document that retail investors appear to ``chase after big orders'' that they suspect are from NT. This pattern is very NT-specific and is not simply a general tendency to chase big orders. We show this by splitting 2015 into an ``NT-aware'' regime versus an ``NT-unaware'' regime. In the former, they chased; in the latter, they did not.

% Beliefs
We then go one step further. The model shows that the magnitude of indirect effects arises from investor beliefs about future stock prices, which in turn depend on investor beliefs about the scale and commitment of NT support for the market. To gauge investor beliefs, we examine discussions on Guba, focusing on key events in 2015 (and later, if possible). We use large language models to extract investor beliefs about the two parameters in the model: whether NT cares, and whether NT is committed to continued market support. We show that in 2015 it took time for the market to become convinced of these two points; only after that did the intervention become effective. If possible, we should also examine what happened in other episodes (such as 2020 and 2022) to see whether things have changed over time.

% Quantification via multipliers
To quantify the effectiveness of NT intervention, we measure the magnitude of indirect effects as ``intervention multiplier'', defined as follows. Say \$1 of direct purchase can only create a current price effect of \$x via demand effects without announcement. If the accompanied NT announcement made price move by $\$x \times y$, we say that the multiplier is $y$, which (roughly) means that ``it is as if NT did $y$ times buying''. $y=1$ means that the indirect effects are entirely useless. $y>1$ means that there is a ``multiplier'' effect, which means that NT's intervention is more effective (per unit of capital committed or risk taken). 

% Final takeaway
Our central insight is, for intervention to be effective/efficient --- measured by a high ``intervention multiplier'' --- such intervention needs to be disclosed, and credibly committed. Ideally we would like to be able to show something that points at this direction. It  might be worth comparing the interventions in 2015 and later during COVID. It seems that the latest estimate of intervention multipliers are higher during the COVID period. I wonder if it can be attributed to our story. (Of course, there are also a million other differences). 

% Possible objections to our policy recommendation
\textcolor{blue}{J: subsequent discussion are less important... they are more specific to China}
If disclosure/commitment is so important, why did NT in 2015 initially shied away from doing this? There are at least a few possibilities, some of which we consider reasonable (though cannot be studied within standard economics framework), while some others we are going to dispell. 1) There is a concern about NT being front-run. At least when it comes to disclosure of NT trading ex-post, this is entirely misguided. Even for ex-ante forward guidance, as long as the guidance is of some kind of ``conditional intervention'' form, it is easy to show that there is no such concern. If NT profit and loss is the concern, the announcement effect further increases NT profits. This is something we can show in the data. In terms of policy effectiveness, ``front-running'' is exactly what NT wants --- ideally there is so much front-running that it doesn't even need to go in and buy. 

Another possible concern about NT disclosure is that NT may have limited ammunition---that is, a position-size constraint or a maximum-loss constraint. The idea is that if NT cannot suffer losses above some threshold, then it does not want others to know it is close to that threshold, because this might invite predatory trading. This can be a valid concern and requires modifying our theory to introduce strategic players.


\section{Model}
\label{sec:model}
A simple two-period model.

\section{Empirical Evidence about Announcement Effects}
\label{sec:announcements}
Document the announcement effects of NT intervention. Establish causality: NT intervention impacts prices and investor demand. 

\section{Empirical Evidence about Beliefs}
\label{sec:beliefs}
Look at Guba posts to study how investor beliefs respond to NT intervention. 

\section{Intervention Multiplier}
\label{sec:multiplier}
Measure the intervention multiplier. Trying to quantify the effectiveness of NT intervention, basically. 

\section{Conclusion}
\label{sec:conclusion}


\clearpage
\bibliographystyle{rfs}
\bibliography{refs}

\processdelayedfloats

\newpage
\appendix

    
 % Resetting for the actual figures and tables
\setcounter{figure}{0}
\setcounter{table}{0}
\renewcommand{\thetable}{\Alph{section}\arabic{table}}
\renewcommand{\thefigure}{\Alph{section}\arabic{figure}}

% \section{Additional Results}
% \label{app:additional}


\processdelayedfloats

\end{document}
